\SourceLeaf{224}
\fancyhead[L]{\small DISQUISITIONES GENERALES CIRCA SUPERFICIES CURVAS}
\renewcommand{\Article}[1]{\par\medskip\Needspace{4\baselineskip}\pdfbookmark[1]{Art. #1}{superficies-art-#1}\centerline{\textbf{#1.}}\nobreak\smallskip}
\begin{center}
DISQUISITIONES GENERALES\par\medskip
{\large CIRCA SUPERFICIES CURVAS.}\par\medskip
\rule{.65in}{.3pt}
\end{center}
\Article{1}
\indent Disquisitiones, in quibus de directionibus variarum rectarum in spatio agitur. plerumque ad maius perspicuitatis et simplicitatis fastigium evehuntur, in auxilium vocando superficiem sphaericam radio \(=1\) circa centrum arbitrarium descriptam, cuius singula puncta repraesentare censebuntur directiones rectarum radiis ad illa terminatis parallelarum. Dum situs omnium punctorum in spatio per tres coordinatas determinatur, puta per distantias a tribus planis fixis inter se normalibus, ante omnia considerandae veniunt directiones axium his planis normalium: puncta superficiei sphaericae, quae has directiones repraesentant, per \((1),(2),(3)\) denotabimus; mutua igitur horum distantia erit quadrans. Ceterum axium directiones versus eas partes acceptas supponemus, versus quas coordinatae respondentes crescunt.

\Article{2}
\indent Haud inutile erit, quasdam propositiones, quae in huiusmodi quaestionibus usum frequentem offerunt, hic in conspectum producere.

I. Angulus inter duas rectas se secantes mensuratur per arcum inter puncta, quae in superficie sphaerica illarum directionibus respondent.

II. Situs cuiuslibet plani repraesentari potest per circulum maximum in superficie sphaerica, cuius planum illi est parallelum.
